% ID: SEL-D0034-OA-E03
% ESTADO: cotejada. Segunda pasada visual y borrador compilado revisados.
% Fuente: Selectividad/matematicas_aplicadas/titular_septiembre_1_MATEMATICAS_APLICADAS.pdf
% SHA-256: c637bbde5c26c0b35b59202bd832ddc97990d3845365ce483fbf2a49832f73fe
% Página: 1; opción A; ejercicio original 3.

Sean $A$ y $B$ dos sucesos aleatorios tales que
\[
P(A)=0.3,\qquad P(B)=0.6,\qquad P(A^C\cap B^C)=0.28.
\]
\begin{enumerate}[label=\alph*)]
\item \textbf{(1 punto)} Halle la probabilidad de que ocurran ambos sucesos a la vez.
\item \textbf{(1 punto)} Calcule la probabilidad de que ocurra $A$ sabiendo que no ha ocurrido $B$.
\item \textbf{(0.5 puntos)} \textquestiondown{}Son $A$ y $B$ independientes?
\end{enumerate}
